% ===== FILE 2/9: Pinhole, EWA projection, Compact box, Trường alpha, Alpha blending, Transmittance decay =====

% --- Slide: Chiếu camera pinhole ---
\begin{frame}{Chiếu điểm 3D lên màn hình: mô hình pinhole}
\scriptsize
Để vẽ một Gaussian 3D lên ảnh 2D, bước đầu tiên là biết \textbf{tâm của nó rơi vào pixel nào}. Camera pinhole (mô hình
camera tiêu chuẩn trong thị giác máy tính) chiếu một điểm $(x,y,z)$ trong hệ toạ độ camera thành pixel $(u,v)$ theo
phép chiếu phối cảnh cổ điển:
\[
\boxed{\;u=f_x\frac{x}{z}+c_x,\qquad v=f_y\frac{y}{z}+c_y\;}
\]
với $(f_x,f_y)$ là tiêu cự tính theo đơn vị pixel (phụ thuộc độ phân giải ảnh và góc nhìn của camera), $(c_x,c_y)$ là
\textbf{điểm chính} (thường là chính giữa ảnh, $W/2,H/2$).
\begin{itemize}\itemsep1pt
    \item Đây là phép chiếu \textbf{phi tuyến} vì có phép chia cho $z$ (độ sâu) --- vật càng xa càng nhỏ lại, đúng như
    cảm nhận thị giác thông thường.
    \item Vấn đề: Gaussian không phải một điểm mà là cả một \textbf{ellipsoid 3D} --- phép chiếu phi tuyến của cả một
    khối hình không cho ra một ellipse đẹp một cách chính xác. Giải pháp thực dụng: \textbf{tuyến tính hoá cục bộ} quanh
    tâm Gaussian bằng ma trận Jacobian $J$ (đạo hàm riêng của phép chiếu tại đúng điểm đó):
\end{itemize}
\[
J=\begin{pmatrix} f_x/z & 0 & -f_x x/z^2\\ 0 & f_y/z & -f_y y/z^2\end{pmatrix}
\]
Ma trận $J$ này sẽ được dùng lại ở \textbf{cả hai nơi}: chiếu hiệp phương sai sang 2D ngay sau đây, và chiếu trục
Gaussian của cơ chế SADGS về sau.
\begin{center}
\includegraphics[width=0.34\linewidth]{ch3_pinhole_projection.png}
\end{center}
\end{frame}

% --- Slide: EWA covariance projection ---
\begin{frame}{EWA splatting: chiếu ellipsoid 3D thành ellipse 2D}
\scriptsize
Dùng Jacobian $J$ ở slide trước, ta xấp xỉ tuyến tính phép chiếu của toàn bộ hiệp phương sai $\Sigma$ (kỹ thuật
\emph{Elliptical Weighted Average}, Zwicker et al. 2001 --- nền tảng toán học cho mọi hệ splatting hiện đại):
\[
\boxed{\;\Sigma' = \big[\,J\,W\,\Sigma\,W^\top J^\top\,\big]_{2\times2} + \kappa\, I\;}
\]
$W$ là ma trận quay world$\to$camera; $[\cdot]_{2\times2}$ nghĩa là chỉ giữ lại khối $2{\times}2$ trên-trái của kết quả
(bỏ chiều depth, vì ta chỉ cần ellipse trên mặt phẳng ảnh).
\begin{itemize}\itemsep1pt
    \item \textbf{Vì sao cần $+\kappa I$?} Khi Gaussian ở rất xa camera hoặc bị nhìn gần như từ cạnh, ellipse chiếu ra
    có thể suy biến thành một đường gần như không có bề rộng --- gây chia cho $0$ và alias (răng cưa). Cộng thêm
    $\kappa I$ đảm bảo ellipse luôn rộng \textbf{ít nhất khoảng 1 pixel} theo mọi hướng --- một bộ lọc low-pass đơn giản, \textbf{luôn được áp dụng}, không phải tuỳ chọn.
    \item Trong CUDA kernel thực tế của repo này (\texttt{cuda\_rasterizer/forward.cu:113--117}), $\kappa=0.1$ --- khác
    với con số $0.3$ vẫn thường thấy trong một số tài liệu 3DGS khác.
    \item Vì "phồng" hiệp phương sai làm sai lệch tổng năng lượng của Gaussian, renderer nhân bù ngược một \textbf{hệ số
    hiệu chỉnh} $\mathrm{coef}=\sqrt{\det\Sigma'_{\text{trước}}/\det\Sigma'_{\text{sau}}}$ vào $\alpha$ để giữ năng lượng gần đúng.
\end{itemize}
\begin{center}
\includegraphics[width=0.34\linewidth]{ch3_ewa_projection.png}
\end{center}
\end{frame}

% --- Slide: Compact box ---
\begin{frame}{Compact box: xác định Gaussian phủ những tile nào}
\scriptsize
Sau khi có ellipse 2D, renderer cần biết nó \textbf{chạm tới những tile $16\times16$ pixel nào} trên màn hình --- đây
là bước quyết định renderer nhanh hay chậm (Gaussian chạm càng nhiều tile, càng tốn công xử lý). Repo này dùng kỹ thuật
\textbf{SnugBox} (\texttt{cuda\_rasterizer/auxiliary.h:313--340}, kỹ thuật kế thừa từ dòng nghiên cứu Taming/Speedy-Splat,
\textbf{không phải điểm mới của SADGS}) --- thay vì bao Gaussian bằng một hình vuông cố định $3\sigma$ như 3DGS gốc,
SnugBox dựng một hộp bao \textbf{khít sát theo đúng hình dạng ellipse thật}, tại ngưỡng nơi $\alpha\cdot G=1/255$ (dưới
mức này coi như không nhìn thấy được nữa), có thể siết chặt hơn nữa bằng một hệ số điều chỉnh \texttt{mult}:
\[
\boxed{\;t=\texttt{mult}\cdot 2\ln(255\,\alpha)\;}\qquad(\texttt{mult}=0.7\text{ mặc định}, \texttt{arguments/\_\_init\_\_.py:114})
\]
\begin{itemize}\itemsep1pt
    \item $\texttt{mult}=1$: khớp đúng ngưỡng $1/255$ gốc (SnugBox nguyên bản, không cắt xén thêm gì).
    \item $\texttt{mult}<1$: cắt \textbf{sớm hơn} $1/255$ --- ngưỡng cắt thật trở thành $\alpha_{\text{biên}}=\sigma_o(255\sigma_o)^{-\texttt{mult}}>1/255$, nghĩa là đánh đổi một phần đuôi Gaussian \emph{vẫn còn nhìn thấy được} để lấy tốc độ (ít tile hơn/Gaussian $\Rightarrow$ ít công tính hơn).
\end{itemize}
\begin{center}
\includegraphics[width=0.30\linewidth]{ch3_compact_box.png}
\end{center}
\end{frame}

% --- Slide: Trường alpha ---
\begin{frame}{Trường alpha: mỗi Gaussian che phủ pixel bao nhiêu?}
\scriptsize
Một Gaussian không "đặc" đều --- độ che phủ giảm dần từ tâm ra biên theo đúng hàm $G(x)$. Tại một pixel cụ thể $x$
trên màn hình, độ mờ đục hiệu dụng của Gaussian thứ $n$ là:
\[
\boxed{\;\alpha_n(x)=\min\big(0.99,\ \alpha_n\, G_n(x)\big)\;}
\]
\begin{itemize}\itemsep1pt
    \item $\alpha_n$ (tham số học được, cố định cho cả Gaussian) nhân với $G_n(x)$ (giá trị hàm Gaussian 2D đã chiếu,
    \textbf{thay đổi theo từng pixel} $x$ --- cao nhất tại tâm ellipse, giảm dần ra ngoài).
    \item Chặn trên $0.99$ (không phải $1.0$): tránh một Gaussian duy nhất "che kín hoàn toàn" khiến gradient của mọi
    Gaussian phía sau nó biến mất hẳn --- luôn để lại một khe hở rất nhỏ cho tín hiệu gradient lan truyền.
    \item Đây chính là đại lượng $\alpha_n(x)$ sẽ được dùng trong công thức alpha blending ở slide tiếp theo --- lưu ý nó
    phụ thuộc pixel, khác với $\alpha_n$ hằng số của riêng Gaussian.
\end{itemize}
\begin{center}
\includegraphics[width=0.36\linewidth]{ch4_alpha_field.png}
\end{center}
\end{frame}

% --- Slide: Alpha blending ---
\begin{frame}{Alpha blending: trộn nhiều lớp Gaussian theo thứ tự trước-sau}
\scriptsize
Tại một pixel, có thể có \textbf{nhiều Gaussian} chồng lên nhau theo chiều sâu. Màu cuối cùng là kết quả trộn tuần tự
từ Gaussian gần camera nhất ra xa nhất (front-to-back alpha blending, giống cách vẽ các lớp trong suốt chồng lên nhau):
\[
\boxed{\;C(x)=\sum_{n} c_n\,\alpha_n(x)\,T_n(x) \;+\; T_{\text{final}}(x)\,C_{bg}\;},\qquad
T_n(x)=\prod_{j<n}\big(1-\alpha_j(x)\big)
\]
\begin{itemize}\itemsep1pt
    \item $T_n$ (\textbf{transmittance}): tỉ lệ ánh sáng \emph{còn lại}, chưa bị các Gaussian phía trước "hấp thụ" ---
    Gaussian thứ $n$ chỉ đóng góp đúng phần ánh sáng nó nhận được sau khi đã đi qua $n-1$ lớp trước nó.
    \item Số hạng $T_{\text{final}}C_{bg}$ (màu \textbf{nền}) là bắt buộc, không được bỏ qua: nếu tổng $\alpha_n(x)$ dọc
    một tia không phủ hết $1$ (ví dụ vùng thưa Gaussian), phần còn thiếu phải được lấp bằng màu nền, nếu không ảnh sẽ bị
    tối sai so với thực tế.
\end{itemize}
\begin{center}
\includegraphics[width=0.36\linewidth]{ch4_alpha_blend_stack.png}
\end{center}
\end{frame}

% --- Slide: Transmittance decay & early stop ---
\begin{frame}{Transmittance suy giảm: cơ sở cho việc dừng sớm}
\scriptsize
Vì $T_n=\prod_{j<n}(1-\alpha_j(x))$ là tích của các thừa số $\le1$, $T_n$ \textbf{giảm dần theo cấp số nhân} khi đi qua
ngày càng nhiều lớp Gaussian (đặc biệt rõ khi từng $\alpha_j(x)$ đều đáng kể). Sau một số lớp đủ nhiều, $T_n$ trở nên
cực nhỏ --- các Gaussian tiếp theo dù có màu gì cũng gần như \textbf{không còn ảnh hưởng} tới $C(x)$ nữa vì bị nhân với
một hệ số gần $0$.
\begin{itemize}\itemsep1pt
    \item Renderer tận dụng điều này để \textbf{dừng sớm} (early termination): ngay khi $T\cdot(1-\alpha)<10^{-4}$ tại
    một pixel, ngừng xử lý thêm Gaussian cho pixel đó --- tiết kiệm đáng kể thời gian tính toán mà hầu như không ảnh
    hưởng chất lượng ảnh (sai số dưới ngưỡng nhận biết được).
    \item Đây cũng là lý do vì sao \textbf{thứ tự depth} (Gaussian nào gần camera hơn) rất quan trọng: nếu sắp sai thứ
    tự, dừng sớm có thể bỏ lỡ một Gaussian đáng lẽ vẫn còn đóng góp đáng kể.
\end{itemize}
\begin{center}
\includegraphics[width=0.36\linewidth]{ch4_transmittance_decay.png}
\end{center}
\end{frame}
